\section{Model interpretability}\label{sec:interp}
No SHAP analysis was run; the notebooks use model-agnostic permutation importance (increase in validation mean squared error when a feature is shuffled), grouped importance by feature family, Ridge coefficients and importance-stability diagnostics, which is consistent with the known pitfalls of such methods \citep{molnar2020interpretability,fisher2019mcr}. Table~\ref{tab:importance} lists the most influential inputs of the target-free Ridge benchmark model.

\begin{table}[!htbp]\centering\small\zebra
\caption{Largest permutation importances, target-free benchmark on the full feature set.}\label{tab:importance}
\begin{tabularx}{\textwidth}{L{0.08\textwidth}YL{0.20\textwidth}R{0.17\textwidth}}\toprule
\thead\hd{Rank} & \hd{Feature} & \hd{Family} & \hd{MSE increase}\\\midrule
1 & \var{FLOW\_RATE\_PHOSPHORIC\_ACID\_2} & raw & 0.805\\
2 & \var{FLOW\_RATE\_PHOSPHORIC\_ACID\_1} & raw & 0.401\\
3 & \var{PRESSURE\_VAPEUR} & raw & 0.150\\
4 & \var{STEAM\_ENERGY\_PROXY} & process-aware & 0.127\\
5 & \var{RECYCLAGE} & raw & 0.082\\\bottomrule
\end{tabularx}
\end{table}

By family, raw sensors account for most of the importance (1.77), followed by temporal features (0.28) and process-aware features (0.13); missingness indicators contribute nothing. The two acid flows dominate, which agrees with the chemistry of an acid attack on rock, but the result describes predictive association. It identifies which signals carry information about the free acid and is used to check that the model relies on plausible process drivers (acid feeds, steam and recycle), not to recommend changing a variable. The same caution applies to the multivariate normalised process-change score used in the notebooks and dashboard: it measures joint sensor movement, not target change, a confirmed startup or shutdown, or a model output.
